% from ega2-1-fr.tex
\oldpage[II]{100}

  \item[\rm{(iii)}] Si $f:X\to Y$ est un $S$-morphisme quasi-projectif,
  $f_{(S')}:X_{(S')}\to Y_{(S')}$ est quasi-projectif pour toute extension
  $S'\to S$ du préschéma de base.
  \item[\rm{(iv)}] Si $f:X\to Y$ et $g:X'\to Y'$ sont deux $S$-morphismes
  quasi-projectifs, $f\times_Sg$ est quasi-projectif.
  \item[\rm{(v)}] Si $f:X\to Y$, $g:Y\to Z$ sont deux morphismes tels que
  $g\circ f$ soit quasi-projectif, et si $g$ est séparé ou $X$ localement
  noethérien, alors $f$ est quasi-projectif.
  \item[\rm{(vi)}] Si $f$ est un morphisme quasi-projectif, il en est de
  même de $f_{\mathrm{red}}$.
\end{enumerate}
\end{proposition}

(i) résulte de (\hyperref[II.5.1.6-fr]{5.1.6}) et
(\hyperref[II.5.1.10-fr]{5.1.10, (i)}). Les autres assertions sont des
conséquences immédiates de la définition (\hyperref[II.5.3.1-fr]{5.3.1}),
des propriétés des morphismes de type fini (\hyperref[I.6.3.4-fr]{I, 6.3.4})
et de (\hyperref[II.4.6.13-fr]{4.6.13}).

\begin{remark}[5.3.5]
\label{II.5.3.5-fr}
On notera qu'il peut se faire que $f_{\mathrm{red}}$ soit quasi-projectif
sans que $f$ le soit, même en supposant que $Y$ est le spectre d'une algèbre
de rang fini sur $\mathbf{C}$ et que $f$ est propre.
\end{remark}

\begin{corollary}[5.3.6]
\label{II.5.3.6-fr}
Si $X$ et $X'$ sont deux $Y$-schémas quasi-projectifs, $X\amalg X'$ est un
$Y$-schéma quasi-projectif.
\end{corollary}

Cela résulte de (\hyperref[II.4.6.18-fr]{4.6.18}).

\subsection{Morphismes propres et morphismes universellement fermés.}
\label{subsection:II.5.4-fr}
\label{II.5.4-fr}

\begin{definition}[5.4.1]
\label{II.5.4.1-fr}
On dit qu'un morphisme de préschémas $f:X\to Y$ est \emph{propre} s'il
vérifie les deux conditions suivantes~:
\begin{enumerate}
  \item[\rm{a)}] $f$ est séparé et de type fini.
  \item[\rm{b)}] Pour tout préschéma $Y'$ et tout morphisme $Y'\to Y$, la
  projection $f_{(Y')}:X\times_YY'\to Y'$ est un morphisme fermé
  (\hyperref[I.2.2.6-fr]{I, 2.2.6}).
\end{enumerate}
Lorsqu'il en est ainsi, on dit aussi que $X$ (considéré comme $Y$-préschéma
de morphisme structural $f$) est propre au-dessus de $Y$.
\end{definition}

Il est immédiat que les conditions a) et b) sont locales sur $Y$. Pour
vérifier que l'image d'une partie fermée $Z$ de $X\times_YY'$ par la
projection $q:X\times_YY'\to Y'$ est fermée dans $Y'$, il suffit de voir que
$q(Z)\cap U'$ est fermé dans $U'$ pour tout ouvert affine $U'$ de $Y'$~;
comme $q(Z)\cap U'=q(Z\cap q^{-1}(U'))$ et que $q^{-1}(U')$ s'identifie à
$X\times_YU'$ (\hyperref[I.4.4.1-fr]{I, 4.4.1}), on voit que pour vérifier
la condition b) de la déf. (\hyperref[II.5.4.1-fr]{5.4.1}), on peut se
limiter au cas où $Y'$ est un schéma affine. On verra plus loin
(\hyperref[II.5.6.3-fr]{5.6.3}) que si $Y$ est localement noethérien, on
peut même se borner à vérifier b) lorsque $Y'$ est de type fini sur $Y$.

Il est clair que tout morphisme propre est \emph{fermé}.

\begin{proposition}[5.4.2]
\label{II.5.4.2-fr}
\begin{enumerate}
  \item[\rm{(i)}] Une immersion fermée est un morphisme propre.
  \item[\rm{(ii)}] Le composé de deux morphismes propres est propre.
  \item[\rm{(iii)}] Si $X$, $Y$ sont des $S$-préschémas,
  $f:X\to Y$ un $S$-morphisme propre, alors
  $f_{(S')}:X_{(S')}\to Y_{(S')}$ est propre pour toute extension
  $S'\to S$ du préschéma de base.
  \item[\rm{(iv)}] Si $f:X\to Y$ et $g:X'\to Y'$ sont deux $S$-morphismes
  propres, le $S$-morphisme
  $f\times_Sg:X\times_SX'\to Y\times_SY'$ est propre.
\end{enumerate}
\end{proposition}

